\chapter[Desenvolvimento]{Desenvolvimento}
\label{cap:desenvolvimento}

O objetivo deste trabalho é o desenvolvimento de uma aplicação web que possibilite o contato e agendamento de sessões entre tutores e aprendizes que precisem de tutoria em uma ou mais áreas de conhecimento, contando com um sistema de recomendação de tutores à aprendizes, empregando elementos de gamificação para manter os usuários mais engajados no uso do software proposto por este trabalho.

A pesquisa bibliográfica realizada permitiu o embasamento teórico e o melhor
entendimento do contexto do projeto: a aprendizagem colaborativa por tutoria, o potencial dos sistemas de recomendação híbridos e a aplicação de gamificação para o engajamento de usuários. A definição da metodologia de pesquisa (pesquisa aplicada e exploratória) permitiu a estruturação das necessidades e a elicitação dos requisitos por meio de métodos de identificação dos personas e as histórias de usuário. Com isso, foi possível delimitar o escopo e definir as metodologias de desenvolvimento ágil (XP e \textit{Kanban}) que guiarão a implementação de toda a proposta apresentada neste trabalho (TCC1).

Para a validação da arquitetura proposta, foram parcialmente implementadas três histórias de usuário em uma simulação simples. Diante do sucesso da simulação, foi demonstrada a viabilidade da execução do projeto para as demais funcionalidades a serem desenvolvidas no decorrer da próxima etapa do projeto (TCC2).

Portanto, a proposta de software apresentada neste trabalho poderá auxiliar aprendizes a encontrarem tutores para lhes auxiliarem em seus aprendizados, além de possibilitarem aos tutores expandirem sua clientela ou público-alvo, contribuindo para mitigação de problemas de aprendizagem de estudantes brasileiros em diferentes áreas de conhecimento envolvendo interesses pessoais ou profissionais.